\documentclass[11pt, a4paper]{article}
\usepackage[a4paper, top=1.8cm, bottom=1.8cm, left=2.0cm, right=2.0cm]{geometry}

\usepackage[hungarian]{babel}


\usepackage{enumitem}
\setlist[itemize]{label=-, leftmargin=1.6em, itemsep=1.5pt, topsep=3pt, parsep=0pt}
\usepackage{amsmath, amssymb}
\usepackage{xcolor}

%%% Megjelenítéshez automatikusan kiegészített csomagok %%%
\usepackage[utf8]{inputenc}
\usepackage[T1]{fontenc}
\usepackage{amsfonts, amsthm}
\usepackage{tcolorbox}
\tcbuselibrary{theorems,skins,breakable}
\usepackage{tikz}
\usepackage{pgfplots}
\pgfplotsset{compat=1.18}
\usepackage{graphicx}
\usepackage{booktabs}
\usepackage{float}
\usepackage{caption}
\usepackage{microtype}

\usepackage[hidelinks]{hyperref}

\definecolor{primary}{HTML}{1E3A5F}
\definecolor{muted}{HTML}{4A5568}
\definecolor{rulegray}{HTML}{CBD5E0}

\setlength{\parindent}{0pt}
\setlength{\parskip}{5pt}

\begin{document}
\setlength{\abovedisplayskip}{6pt}
\setlength{\belowdisplayskip}{6pt}
\setlength{\abovedisplayshortskip}{4pt}
\setlength{\belowdisplayshortskip}{4pt}

% --- Fejléc ---
{\color{primary}\Large\bfseries 3D Koordinátageometria – Dolgozat (A)}\hfill{\color{muted}\large\bfseries Max. 24 pont}\\[3pt]
{\color{muted}\normalsize Kidolgozott javítókulcs és részletes pontozási útmutató}
\vspace{4pt}
{\color{primary}\hrule height 1.2pt}
\vspace{10pt}

% --- 1. feladat ---
{\color{primary}\bfseries 1. feladat (2 pont)}

\textit{Adott két pont a térben: $A(1,2,3)$ és $B(1,6,6)$. Számítsd ki a két pont távolságát!}

\textbf{Megoldás:} Két térbeli pont távolságának képlete alapján:
\[
\begin{aligned}
d(A,B) &= \sqrt{(x_B - x_A)^2 + (y_B - y_A)^2 + (z_B - z_A)^2} \\
       &= \sqrt{(1 - 1)^2 + (6 - 2)^2 + (6 - 3)^2} = \sqrt{0 + 16 + 9} = \sqrt{25} = 5
\end{aligned}
\]
\textbf{Válasz:} A két pont távolsága \textbf{5} egység.

{\color{muted}\small\textbf{Pontozás (2 pont):}}
\begin{itemize}
    \item {\color{muted}\small A távolságképlet helyes felírása és behelyettesítése: \textbf{1 pont}}
    \item {\color{muted}\small Helyes számolás és végeredmény ($5$): \textbf{1 pont}}
\end{itemize}

\vspace{6pt}{\color{rulegray}\hrule}\vspace{8pt}

% --- 2. feladat ---
{\color{primary}\bfseries 2. feladat (2 pont)}

\textit{Számítsd ki az $\mathbf{a}(2,1,-2)$ és $\mathbf{b}(3,0,4)$ vektorok skaláris szorzatát!}

\textbf{Megoldás:} A skaláris szorzat koordinátás képletével:
\[
\mathbf{a} \cdot \mathbf{b} = a_1 b_1 + a_2 b_2 + a_3 b_3 = 2 \cdot 3 + 1 \cdot 0 + (-2) \cdot 4 = 6 + 0 - 8 = -2
\]
\textbf{Válasz:} A két vektor skaláris szorzata \textbf{$-2$}.

{\color{muted}\small\textbf{Pontozás (2 pont):}}
\begin{itemize}
    \item {\color{muted}\small A skaláris szorzat képletének alkalmazása és behelyettesítés: \textbf{1 pont}}
    \item {\color{muted}\small Helyes előjelkezelés és végeredmény ($-2$): \textbf{1 pont}}
\end{itemize}

\vspace{6pt}{\color{rulegray}\hrule}\vspace{8pt}

% --- 3. feladat ---
{\color{primary}\bfseries 3. feladat (3 pont)}

\textit{Írd fel a $P_0(2,-1,4)$ ponton átmenő, $\mathbf{v}(3,5,-1)$ irányvektorú egyenes paraméteres egyenletrendszerét!}

\textbf{Megoldás:} A $P_0(x_0, y_0, z_0)$ ponton átmenő, $\mathbf{v}(v_1, v_2, v_3)$ irányvektorú egyenes paraméteres egyenletrendszere:
\[
\begin{cases}
x = x_0 + t \cdot v_1 \\
y = y_0 + t \cdot v_2 \\
z = z_0 + t \cdot v_3
\end{cases}
\implies
\begin{cases}
x = 2 + 3t \\
y = -1 + 5t \\
z = 4 - t
\end{cases}
\qquad (t \in \mathbb{R})
\]
{\color{muted}\small\textbf{Pontozás (3 pont):}}
\begin{itemize}
    \item {\color{muted}\small Az $x = 2 + 3t$ egyenlet helyes felírása: \textbf{1 pont}}
    \item {\color{muted}\small Az $y = -1 + 5t$ egyenlet helyes felírása: \textbf{1 pont}}
    \item {\color{muted}\small A $z = 4 - t$ egyenlet helyes felírása: \textbf{1 pont}}
\end{itemize}

\vspace{6pt}{\color{rulegray}\hrule}\vspace{8pt}

% --- 4. feladat ---
{\color{primary}\bfseries 4. feladat (3 pont)}

\textit{Írd fel a $P_0(1,1,1)$ ponton átmenő, $\mathbf{n}(2,-1,3)$ normálvektorú sík egyenletét! Hozza az egyenletet $Ax+By+Cz=D$ alakra!}

\textbf{Megoldás:} A sík normálvektoros egyenlete $A(x - x_0) + B(y - y_0) + C(z - z_0) = 0$, amelybe behelyettesítve:
\[
\begin{aligned}
2(x - 1) - 1(y - 1) + 3(z - 1) &= 0 \\
2x - 2 - y + 1 + 3z - 3 &= 0 \\
2x - y + 3z &= 4
\end{aligned}
\]
\textit{(Alternatív megoldás: $2x - y + 3z = 2\cdot 1 + (-1)\cdot 1 + 3\cdot 1 = 4$.)}

\textbf{Válasz:} A sík egyenlete \textbf{$2x - y + 3z = 4$}.

{\color{muted}\small\textbf{Pontozás (3 pont):}}
\begin{itemize}
    \item {\color{muted}\small A normálvektor és a pont helyes behelyettesítése: \textbf{1 pont}}
    \item {\color{muted}\small Zárójelek felbontása és helyes összevonás: \textbf{1 pont}}
    \item {\color{muted}\small A kért $Ax + By + Cz = D$ alak megadása ($2x - y + 3z = 4$): \textbf{1 pont}}
\end{itemize}

\newpage

% --- 5. feladat ---
{\color{primary}\bfseries 5. feladat (2 pont)}

\textit{Adja meg az $O(-2,3,0)$ középpontú, $r=4$ sugarú gömb egyenletét!}

\textbf{Megoldás:} Az $O(u, v, w)$ középpontú, $r$ sugarú gömb egyenlete $(x - u)^2 + (y - v)^2 + (z - w)^2 = r^2$:
\[
(x - (-2))^2 + (y - 3)^2 + (z - 0)^2 = 4^2 \implies (x + 2)^2 + (y - 3)^2 + z^2 = 16
\]
\textbf{Válasz:} A gömb egyenlete \textbf{$(x + 2)^2 + (y - 3)^2 + z^2 = 16$}.

{\color{muted}\small\textbf{Pontozás (2 pont):}}
\begin{itemize}
    \item {\color{muted}\small A bal oldal előjelhelyes felírása $(x+2)^2 + (y-3)^2 + z^2$: \textbf{1 pont}}
    \item {\color{muted}\small A jobb oldal helyes értéke ($r^2 = 16$): \textbf{1 pont}}
\end{itemize}

\vspace{7pt}{\color{rulegray}\hrule}\vspace{9pt}

% --- 6. feladat ---
{\color{primary}\bfseries 6. feladat (8 pont)}

\textit{Határozza meg az $e:\dfrac{x-1}{2}=\dfrac{y-2}{3}=\dfrac{z-3}{-1}$ egyenletrendszerrel adott egyenes és az $S:2x-y+3z=11$ sík döféspontját (azaz metszéspontjának koordinátáit)!}

\textbf{Megoldás:}

\textbf{1. lépés – Az egyenes felírása paraméteres alakban:} Legyen a közös arány $t \in \mathbb{R}$:
\[
\frac{x-1}{2} = \frac{y-2}{3} = \frac{z-3}{-1} = t \implies x = 1 + 2t, \quad y = 2 + 3t, \quad z = 3 - t
\]

\textbf{2. lépés – Behelyettesítés a sík egyenletébe ($2x - y + 3z = 11$):}
\[
\begin{aligned}
2(1 + 2t) - (2 + 3t) + 3(3 - t) &= 11 \\
2 + 4t - 2 - 3t + 9 - 3t &= 11 \\
9 - 2t &= 11 \\
-2t &= 2 \implies t = -1
\end{aligned}
\]

\textbf{3. lépés – A döféspont ($M$) koordinátáinak meghatározása ($t = -1$ visszahelyettesítésével):}
\[
x = 1 + 2(-1) = -1, \qquad y = 2 + 3(-1) = -1, \qquad z = 3 - (-1) = 4
\]

\textbf{Ellenőrzés:} Az egyenes egyenletében $\dfrac{-1-1}{2} = \dfrac{-1-2}{3} = \dfrac{4-3}{-1} = -1$; a sík egyenletében $2(-1) - (-1) + 3(4) = -2 + 1 + 12 = 11$ (mindkettő teljesül).

\textbf{Válasz:} A döféspont koordinátái: \textbf{$M(-1,\,-1,\,4)$}.

{\color{muted}\small\textbf{Pontozás (8 pont):}}
\begin{itemize}
    \item {\color{muted}\small Az egyenes paraméteres egyenletrendszerének helyes felírása: \textbf{2 pont}}
    \item {\color{muted}\small Behelyettesítés a sík egyenletébe: \textbf{1 pont}}
    \item {\color{muted}\small Zárójelek felbontása és összevonás ($9 - 2t = 11$): \textbf{2 pont}}
    \item {\color{muted}\small A paraméter helyes meghatározása ($t = -1$): \textbf{1 pont}}
    \item {\color{muted}\small A metszéspont mindhárom koordinátájának helyes kiszámítása ($M(-1,-1,4)$): \textbf{2 pont}}
\end{itemize}

\vspace{7pt}{\color{rulegray}\hrule}\vspace{9pt}

% --- 7. feladat ---
{\color{primary}\bfseries 7. feladat (4 pont)}

\textit{Egy ellipszis egyenlete: $\dfrac{x^2}{25}+\dfrac{y^2}{9}=1$. Olvasd le a nagytengely és kistengely felének hosszát, majd számítsd ki a fókuszpont és az origó távolságát!}

\textbf{Megoldás:} Az origó középpontú ellipszis $\dfrac{x^2}{a^2} + \dfrac{y^2}{b^2} = 1$ ($a > b > 0$) kanonikus egyenletéből leolvasható:
\[
a^2 = 25 \implies a = 5 \text{ (nagytengely fele)}, \qquad b^2 = 9 \implies b = 3 \text{ (kistengely fele)}
\]
A fókuszpont és az origó távolsága (a lineáris excentricitás, $c$):
\[
c^2 = a^2 - b^2 = 25 - 9 = 16 \implies c = \sqrt{16} = 4
\]
\textbf{Válasz:} Nagytengely fele \textbf{$a = 5$}, kistengely fele \textbf{$b = 3$}, a fókuszpont és az origó távolsága \textbf{$c = 4$}.

{\color{muted}\small\textbf{Pontozás (4 pont):}}
\begin{itemize}
    \item {\color{muted}\small A nagytengely felének helyes leolvasása ($a = 5$): \textbf{1 pont}}
    \item {\color{muted}\small A kistengely felének helyes leolvasása ($b = 3$): \textbf{1 pont}}
    \item {\color{muted}\small A $c^2 = a^2 - b^2$ összefüggés felírása és behelyettesítése: \textbf{1 pont}}
\item {\color{muted}\small A fókuszpont távolságának helyes kiszámítása ($c = 4$): \textbf{1 pont}}
\end{itemize}

\vfill
{\color{primary}\hrule height 0.8pt}
\vspace{4pt}
{\color{muted}\small\textbf{Pontszámok összesítése:} 1. f. (2p) + 2. f. (2p) + 3. f. (3p) + 4. f. (3p) + 5. f. (2p) + 6. f. (8p) + 7. f. (4p) \hfill \textbf{Összesen: 24 pont}}

\end{document}